\section{Differential ladders as positivity certificates}
\label{sec:ladders}

\subsection{The analytic lemma}

\begin{lemma}[One integrating-factor step]
\label{lem:step}
Let $f,g$ be real functions on $[a,\infty)$, with $f$ continuously
differentiable and $g$ continuous. Suppose
\[
 f'(x)=\rho f(x)+g(x),\qquad f(a)\ge0,\qquad
 g(x)\ge0\quad(x\ge a),
\]
where $\rho\in\R$. Then $f(x)\ge0$ on $[a,\infty)$.
\end{lemma}
\begin{proof}
For $w(x)=e^{-\rho x}f(x)$, the product and chain rules give
$w'(x)=e^{-\rho x}g(x)\ge0$. Hence $w(x)\ge w(a)\ge0$ for $x\ge a$.
Multiplication by the positive number $e^{\rho x}$ gives the conclusion.
Equivalently, the exact identity
\begin{equation}
 f(x)=e^{\rho(x-a)}f(a)+
       \int_a^x e^{\rho(x-t)}g(t)\,dt
 \label{eq:integrating-factor}
\end{equation}
expresses both terms as nonnegative.
\end{proof}

The derivative proof is preferable for a first Lean bridge. It needs only
explicit derivative witnesses, continuity, and monotonicity on a convex ray;
it does not require construction of an integral expression or a separate
integrability proof for every emitted step. Mathlib exposes these ingredients
through \code{HasDerivAt}, exponential derivative lemmas, and the
\code{monotoneOn_of_hasDerivWithinAt_nonneg} theorem~\cite{mathlib-meanvalue,mathlib-expderiv}.

\subsection{The certificate and its soundness}

\begin{definition}
A differential-ladder certificate for $f$ at $a$ consists of real numbers
$\rho_0,\ldots,\rho_{n-1}$ and functions $q_0,\ldots,q_n$ such that
\begin{equation}
 q_0=f,\qquad q_{j+1}=q_j'-\rho_jq_j,\qquad
 q_n=0,\qquad q_j(a)\ge0\quad(0\le j<n).
 \label{eq:ladder}
\end{equation}
In the executable fragment, all $q_j$ are canonical rational exponential
polynomials and all $\rho_j,a$ are rational.
\end{definition}

\begin{theorem}[Ladder soundness]
\label{thm:ladder}
A certificate satisfying \eqref{eq:ladder} proves $f(x)\ge0$ for every $x\ge a$.
\end{theorem}
\begin{proof}
The terminal function is nonnegative. Descend from $q_n$ to $q_0$ using
Lemma~\ref{lem:step}, with $f=q_j$ and $g=q_{j+1}$ at each step.
Every required smoothness condition holds for exponential polynomials.
\end{proof}

The checker does not need to know why the factors were chosen or whether they
are minimal. In particular, it need not trust a spectral calculation. It
checks the derivative identities, the initial inequalities, and the zero
terminal object directly. An unnecessarily long but correct ladder is still a
proof certificate.

\begin{theorem}[Strictness and zero obstruction]
\label{thm:strict}
Suppose a ladder is accepted. If some $q_j(a)>0$, then $f(x)>0$ for every $x>a$.
If $q_0(a)>0$, positivity holds on the closed ray. If every $q_j(a)=0$, then
$f$ is zero throughout the ray. Consequently a nonzero function represented by
an accepted ladder has no zero in $(a,\infty)$.
\end{theorem}
\begin{proof}
If $q_j(a)>0$, equation~\eqref{eq:integrating-factor} shows that $q_j(x)>0$
for $x\ge a$. For any $x>a$, its contribution to the next integral down is
strictly positive. Repeating the argument gives $q_0(x)>0$. When $q_0(a)>0$
the endpoint is positive too. If every initial value is zero, descend from
$q_n=0$ in the same identity to obtain $q_{n-1}=\cdots=q_0=0$.
\end{proof}

This theorem provides more than a strictness flag. It proves that there are
true nonnegativity goals for which \emph{no amount of ladder search} can work.
Section~\ref{sec:enrichment} uses that fact to choose another worker.

\subsection{A certificate which tolerates arbitrarily high boundary contact}

Let
\[
 R_m(x)=e^x-\sum_{k=0}^m\frac{x^k}{k!},\qquad m\ge0.
\]
Take $m+1$ zero factors followed by one factor equal to one. Successive ordinary
derivatives reduce the polynomial degree; after $m+1$ steps the function is
$e^x$, and the final $(\D-1)$ kills it. The initial values are
\[
 \underbrace{0,\ldots,0}_{m+1\text{ values}},1.
\]
Thus $R_m(x)\ge0$ for $x\ge0$ and $R_m(x)>0$ for $x>0$. The certificate grows
linearly in $m$ in its number of steps. Its correctness does not become a
numerical conditioning problem as the contact order increases.

For $m=4$ the exact ladder is
\[
 (0,0,0,0,0,1),\qquad\text{initial vector }(0,0,0,0,0,1).
\]
The delivered experiments exercise $m=0,\ldots,24$. A strict closed-ray request
is deliberately rejected because $R_m(0)=0$; an open-ray request is accepted.

\subsection{The linear-algebra interpretation}

Fix a multiset of $n$ real factors and the space $V$ of solutions of the
associated homogeneous constant-coefficient equation. The map
\[
 f\longmapsto
 \bigl(f(a), (\D-\rho_0)f(a),\ldots,
       (\D-\rho_{n-2})\cdots(\D-\rho_0)f(a)\bigr)
\]
is a coordinate system on $V$. Its coordinates are triangular combinations of
ordinary derivatives at $a$, with diagonal coefficients one. Equivalently,
the first-order triangular system in \eqref{eq:ladder} reconstructs a unique
function from those initial data.

Nonnegative coordinates define a simplicial cone in $V$. Its basis functions
can be constructed backwards from unit initial vectors using
\eqref{eq:integrating-factor}; they are positive on the open ray. Therefore
ladder search is not an unstructured search through arbitrary inequalities.
It is membership testing in a particular positive cone, with a particularly
simple change of coordinates.

This also explains both strengths and weaknesses. The cone is adapted to
boundary zeros, since early coordinates may vanish exactly. But none of its
nonzero elements can have an interior zero, because every nonzero nonnegative
combination of its basis functions is strictly positive there. Adding more
coordinates can enlarge the accepted cone without eliminating that obstruction.

\subsection{A minimal checker loop}

\par\noindent\begin{minipage}{\linewidth}
\begin{lstlisting}
checkLadder(source, anchor, factors, receipt):
    q := source
    strictSeen := false
    for each factor rho and corresponding receipt step:
        require step.polynomial == q
        [lo, hi] := certifiedPointBounds(q, anchor)
        require step.bounds == [lo, hi] and lo >= 0
        strictSeen := strictSeen or (lo > 0)
        q := independentlyComputedShiftedDerivative(q, rho)
    require q == 0
    return NonnegativeRay(anchor, strictSeen)
\end{lstlisting}
\end{minipage}\par

The independently computed derivative in the prototype uses a coefficient
stencil rather than the producer's termwise differentiation loop:
\begin{equation}
 \widetilde c_{\lambda,k}
 = (\lambda-\rho)c_{\lambda,k}+(k+1)c_{\lambda,k+1}.
 \label{eq:stencil}
\end{equation}
This is algorithmic diversity, not a formal verification of either Python
implementation. The separately tested initial-jet oracle uses a third
calculation, described in Section~\ref{sec:experiments}.

A production receipt need not retain every intermediate polynomial: the checker
can reconstruct them from the source and factors. The reference format keeps
them for diagnosis, mutation tests, and readable emitted proofs. Omitting those
fields is a later certificate-size optimisation, not a change in the theorem.
